% Part: first-order-logic
% Chapter: introduction
% Section: substitution

\documentclass[../../../include/open-logic-section]{subfiles}

\begin{document}

\olfileid{fol}{int}{sub}

\olsection{الإحلال}

سنناقش مثالًا يوضح كيف تترابط الأمور، وكيف يضع تطوير النحو والدلالة
أساس بحوثنا الأكثر تقدمًا لاحقًا. فينبغي لأنساق ال!!{derivation} أن
تمكننا من !!{derive} $\Atom{\Obj P}{\Obj a}$ من
$\lforall[\Obj v_0][\Atom{\Obj P}{\Obj v_0}]$. وربما أردنا أن نصوغ
ذلك نفسه قاعدة استدلال. غير أننا، لكي نفعل ذلك، يجب أن نتمكن من صياغته
بأعم صورة: لا لـ~$\Obj P$ و$\Obj a$ و$\Obj v_0$ وحدها، بل لأي
!!{formula}~$!A$، وحد~$t$، و!!{variable}~$x$. (وتذكر أن
ال!!{constant}s حدود، لكننا سنبحث أيضًا حدودًا أعقد مبنية من
ال!!{constant}s وال!!{function}s.) فنريد إذن أن نستطيع قول شيء من
قبيل: «متى فرغت من !!{derive} $\lforall[x][!A(x)]$، جاز لك أن تستنتج
$!A(t)$---أي نتيجة حذف $\lforall[x]$ وإحلال~$t$ محل~$x$». ولكن ما
المقصود بالضبط بـ«إحلال $t$ محل~$x$»؟ وما العلاقة بين $!A(x)$
و~$!A(t)$؟ وهل ينجح ذلك دائمًا؟

ولضبط هذا نعرّف عملية \emph{الإحلال}. والإحلال محفوف بالمزالق في
الحقيقة، إذ لا يجوز أن نستبدل ببساطة~$t$ بكل ورود لـ~$x$ في~$!A$،
ولا يصلح كل~$t$ للإحلال محل كل~$x$. وسنعالج هذا، مرة أخرى، باستعمال
تعريفات استقرائية. لكن متى تم ذلك، أصبحت صياغة قاعدة استدلال على صورة
«استنتج $!A(t)$ من $\lforall[x][!A(x)]$» تعريفًا دقيقًا. وسيمكننا،
فضلًا على ذلك، أن نبين أن هذه قاعدة استدلال جيدة، بمعنى أن
$\lforall[x][!A(x)]$ تستلزم~$!A(t)$. لكن لإثبات ذلك، يلزمنا مرة أخرى
إثبات أمر قد يدفعك لأول وهلة إلى أن تسأل: «لم نفعل هذا؟». فكون
$\lforall[x][!A(x)]$ تستلزم~$!A(t)$ يعتمد على أن حصول
$\Sat{M}{!A(t)}$ أو عدم حصوله لا يتوقف إلا على قيمة الحد~$t$؛ أي إننا
إذا جعلنا $m$ هو ال!!{element} من~$\Domain{M}$ الذي يعيّنه~$t$، كان
$\Sat{M}{!A(t)}[s]$ إذا وفقط إذا كان
$\Sat{M}{!A(x)}[\Subst{s}{m}{x}]$. ويصدق هذا حتى حين يحتوي $t$ على
!!{variable}s، غير أننا سنحتاج إلى العناية بالطريقة الدقيقة التي نصوغ
بها النتيجة.

\end{document}
